\begin{proposition}[Euler number for a bijective normalization]\hypertarget{tag-02YD}{}\label{night:P4R.2}\label{stmt:P4R.2}
\StatementRecord{night:P4R.2}{proved}{02YD}{\SeriesRef{night:P4.12}, \SeriesRef{night:P4.8}, \SeriesRef{night:P4R.1}}

In the standing torus-fibration setting of this section, every irreducible singular support
whose normalization is smooth and whose normalization morphism
is bijective has nonnegative topological Euler number.
Its Euler number is zero precisely when that normalization is
a minimal ruled surface over an elliptic curve or a Hopf surface.
\end{proposition}
\begin{proof}

\begin{enumerate}[label=\arabic*.,leftmargin=*]
\item Let $\nu:T\to S$ denote this finite normalization morphism.
A continuous bijection from the compact $T$ to the Hausdorff
$S$ is a homeomorphism. It therefore induces an isomorphism
of fundamental groups and equality of Euler numbers.
\SeriesRef{night:P4R.1} makes $\pi_1(T)$ virtually abelian.

\item The effective-conductor and canonical-torsion argument of
\SeriesRef{night:P4.12} applies to this same $T$ and proves
$\kappa(T)=-\infty$. Write $q:T\to T_0$ for the contraction
to a minimal model, with $k$ point blowups. The classification
gives rational, ruled, and class-VII possibilities for $T_0$.
In the K\"ahler case, $p_g(T_0)=0$ permits an integral
K\"ahler class; the embedding given by
\cite[Main Theorem, pp.~314--315]{Kodaira1954} identifies $T_0$
with a projective surface, to which
\cite[Theorem~VI.17 and Corollary~VI.18, pp.~111--112]{Beauville1978}
applies. In the non-K\"ahler case use
\cite[Theorem~5.1, p.~27]{OVV2020}. For its elliptic alternative, Theorem~4.2, Remark~4.4 and
Proposition~4.5 there (pp.~17--22) give elliptic reduced fibres,
the cyclic translation model at multiple fibres, and a
holomorphic elliptic action. For the fibration
$\pi:T_0\to C_0$, its nowhere-zero vertical vector field gives
$Q=T_{T_0}/\OO$, and the differential induces
$Q\simeq\pi^*T_{C_0}(-R_\pi)$, where $R_\pi$ is the
multiple-fibre ramification divisor. Thus
$K_{T_0}=\pi^*K_{C_0}(R_\pi)$ and $K_{T_0}^2=0$;
the elliptic fibres also give $e(T_0)=0$. Together with
$p_g=0$, Noether's formula gives $h^1(\OO)=1$.
Theorem~2.5 there (pp.~7--9) then gives $b_1=1$.
This alternative also satisfies the class-VII conditions
$b_1=1$ and $\kappa=-\infty$.
Blowing up a smooth point preserves the fundamental group:
van Kampen compares a four-ball with the disk bundle
of $\OO_{\mathbb P^1}(-1)$, both attached along the same
simply connected boundary $S^3$. For a ruled minimal model
$p:T_0\to C$ over a genus-$g$ curve, the homotopy sequence
of its $\mathbb P^1$-bundle gives the isomorphism $p_*$
on fundamental groups. The resulting isomorphism from
$\pi_1(S)$ to $\pi_1(C)$ is $p_*q_*(\nu_*)^{-1}$.
For $g\ge2$ this group is not virtually abelian. Indeed a
finite-index subgroup corresponds to an unramified cover
of genus $1+d(g-1)>1$, and a genus-$h$ surface group maps
onto the nonabelian free group on $h$ generators by sending
the $b$ generators to the identity in its usual presentation.
Consequently $g\le1$.

\item A point blowup replaces a point by $\mathbb P^1$ and
increases Euler characteristic by one. The cell counts of
$\mathbb P^2$ and a $\mathbb P^1$-bundle over a genus-$g$
curve give $3$ and $4-4g$, respectively. For class VII,
Poincar\'e duality and $b_1(T_0)=1$ give $e(T_0)=b_2(T_0)$.
Thus
\[
 e(T)=\begin{cases}
 3+k,&T_0=\mathbb P^2,\\
 4-4g+k,&T_0\text{ ruled over a genus-}g\text{ curve},\\
 b_2(T_0)+k,&T_0\text{ of class VII}.
 \end{cases}
\]
Since $g\le1$, these formulas give
$e(S)=e(T)\ge0$. Equality in the ruled case is exactly
$g=1$ and $k=0$. Equality in class VII gives a minimal surface with
$b_2=0$. The conductor is nonempty because $S$ is singular
and nonnormal; \SeriesRef{cdp:st:A.5.3} gives its effective
conductor divisor on $T$. By \cite[Theorem~1, p.~1]{BBK2019},
there is a finite \'etale cover $a:T'\to T$ whose domain is
Hopf or Inoue. The inverse image of the conductor contains
a curve, excluding the Inoue alternative by
\cite[Propositions~2(i) and~3(i), pp.~274--275 and~277,
and the paragraph after (21), p.~279]{Inoue1974}.
The composition of $a$ with the Hopf universal cover
$\C^2\setminus\{0\}\to T'$ is a universal cover of $T$;
hence $T$ itself is Hopf. The converse follows from the same Euler
formulas and the homeomorphism $\nu$.

\end{enumerate}
\end{proof}